\section*{Chapitre \ref{ch:mx:transient-impact-and-propagator-models} --- Impact transitoire et modèles de propagateur}

\begin{solution}{exo:mx:transient-impact-and-propagator-models:1}
Avec $\E[\varepsilon_t\varepsilon_{t+1}]=0.27$ et des signes $\pm1$, $\Pr(\text{identiques})-\Pr(\text{différents})=0.27$, donc $\Pr(\text{identiques})=(1+0.27)/2=63.5~\%$.
\end{solution}

\begin{solution}{exo:mx:transient-impact-and-propagator-models:2}
Avec un noyau constant, chaque impact de transaction est permanent, et les signes persistants continuent d'ajouter des impacts de même signe~: le prix dériverait dans le sens du flux,
de façon prévisible, aussi longtemps que dure la mémoire des signes.
\end{solution}

\begin{solution}{exo:mx:transient-impact-and-propagator-models:3}
$\rho T+2=10$~: un bloc de 1\,000 actions au début, 4\,000 actions par heure pendant deux heures, puis un bloc de 1\,000 à la fin.
\end{solution}

\begin{solution}{exo:mx:transient-impact-and-propagator-models:4}
$\Gamma=\begin{pmatrix}1&e^{-1}\\e^{-1}&1\end{pmatrix}$~; l'aller-retour coûte $\tfrac12(1+1-2e^{-1})=1-e^{-1}=0.632$, valeur positive. Pour vendre une unité,
$x\propto\Gamma^{-1}\mathbf 1$, qui est symétrique~: la moitié à chaque instant.
\end{solution}

\begin{solution}{exo:mx:transient-impact-and-propagator-models:5}
Sans mémoire, l'impact de chaque transaction se résorbe et rien ne le remplace en moyenne~: les variations de prix sur de longs retards sont plus petites que la somme des variations courtes, si bien que
la variance par retard diminue (de 0.32 à 0.18).
\end{solution}

\begin{solution}{exo:mx:transient-impact-and-propagator-models:6}
$\alpha-1=0.5$. La pente mesurée, 0.62, est plus raide parce que les ordres des fondamentalistes, dont les signes suivent l'écart à la valeur plutôt que de longues séries, diluent
la mémoire des preneurs de bruit, et que l'ajustement porte sur les retards 2 à 50 d'un échantillon de 10\,764 ordres.
\end{solution}

\begin{solution}{exo:mx:transient-impact-and-propagator-models:7}
Le vecteur propre associé à la valeur propre la plus négative ($-0.181$ sur le sous-espace des allers-retours, soit un profit de 0.091) alterne achats et ventes en sept jambes~: $+0.20$
au départ, $-0.26$ à la quatrième transaction, $+0.33$ à la septième, $-0.37$ à la onzième, et ainsi de suite de façon symétrique. Le noyau reste proche de sa valeur pleine
pendant environ $\tau$, puis chute brusquement~; son spectre est donc négatif à la fréquence de cette oscillation. Chaque jambe échange contre le prix que
sa prédécesseure a laissé avant que cet impact ne se soit résorbé, et une fois qu'il l'a été, selon un schéma qui rapporte davantage qu'il ne coûte.
\end{solution}

\begin{solution}{exo:mx:transient-impact-and-propagator-models:8}
La définition positive exclut les allers-retours profitables, mais non la manipulation déclenchée par les transactions~: un noyau défini positif mais non convexe
($1<\kappa\le2$ dans le chapitre) fait inclure des achats dans la vente la moins coûteuse. La condition sûre est une décroissance convexe et non croissante.
\end{solution}

\begin{solution}{pb:mx:transient-impact-and-propagator-models:1}
\textbf{1.} Signes d'ordres à autocorrélation en loi de puissance, d'exposant environ 0.6 (Hurst 0.7).
\textbf{2.} Découpage de métaordres en loi de puissance en morceaux égaux~: $C(\ell)\sim\ell^{-(\alpha-1)}$.
\textbf{3.} 0.27, 0.106 et 0.023~; un exposant de 0.62.
\textbf{4.} 0.52, 1.05 et 2.3 ticks.
\textbf{5.} $p_t=\sum_{s<t}G(t-s)\varepsilon_s$~; $\E[\varepsilon_{t-k}\Delta p_t]=\sum_nK(n)C(|k-n|)$.
\textbf{6.} 0.541, 0.452 et 0.419 ticks.
\textbf{7.} Marché 1.98 en baisse jusqu'à 1.35~; modèle 0.29 en hausse jusqu'à 1.34~; mélangé 0.32 en baisse jusqu'à 0.18.
\textbf{8.} La décroissance du noyau compense exactement la mémoire des signes, ce qui laisse les prix diffusifs.
\textbf{9.} $\tfrac12x^\top\Gamma x$, $\Gamma_{ij}=G(|t_i-t_j|)$.
\textbf{10.} Minimiser $\tfrac12x^\top\Gamma x-\mu(\mathbf 1^\top x-X)$~: $\Gamma x=\mu\mathbf 1$.
\textbf{11.} Blocs de $X/(\rho T+2)$ aux deux extrémités et un rythme $\rho X/(\rho T+2)$ entre les deux~: le premier bloc utilise le carnet existant, le dernier utilise le carnet
reconstitué, sans transaction ultérieure pour en payer le coût.
\textbf{12.} 1\,000 actions, 4\,000 par heure, 1\,000 actions.
\textbf{13.} Un aller-retour à coût espéré négatif~; une exécution la moins coûteuse comportant des transactions de sens opposés.
\textbf{14.} L'impact doit décroître comme une fonction convexe non croissante du temps.
\textbf{15.} $-0.136$, $-0.042$, $-0.008$, $-0.002$, $+0.028$, $+0.091$, $+0.213$.
\textbf{16.} $\kappa=1.5$ et 1.8 (deux achats chacun), et tout $\kappa>2$.
\textbf{17.} La matrice du noyau gaussien est singulière sur une grille fine~: il s'agit de la frontière de la définie positivité.
\textbf{18.} \emph{Résultat nommé}~: avec $G(t)=e^{-(t/0.3)^\kappa}$ sur 21 transactions, le meilleur aller-retour rapporte 0.028, 0.091 et 0.213 pour $\kappa=2.2$, 2.5 et 3, et
fait perdre pour tout $\kappa\le2$~; la manipulation de type déclenchée par les transactions commence pour $\kappa>1$, où la convexité est perdue.
\textbf{19.} Un plateau suivi d'une chute correspond à $\kappa>1$~: au minimum une manipulation déclenchée par les transactions, et des allers-retours profitables si la chute est brutale ($\kappa>2$).
\textbf{20.} Convexe et non croissante.
\end{solution}

\begin{solution}{iq:mx:transient-impact-and-propagator-models:1}
Chaque impact de transaction se résorbe (impact transitoire), et cette décroissance compense la continuation prévisible du flux d'ordres~; les teneurs de marché qui perçoivent la
persistance du flux ajustent leurs cotations de sorte que les prix restent proches d'une martingale.
\iqlookfor{Impact transitoire~; compensation~; teneurs de marché.}
\end{solution}

\begin{solution}{iq:mx:transient-impact-and-propagator-models:2}
Signes et prix milieux avant transaction~; l'autocorrélation des signes et la corrélation croisée des signes avec les variations ultérieures du prix milieu~; résoudre le système de Toeplitz pour les accroissements du noyau
(Levinson)~; vérifier que le modèle ajusté rend les prix diffusifs.
\iqlookfor{Le système linéaire~; un diagnostic.}
\end{solution}

\begin{solution}{iq:mx:transient-impact-and-propagator-models:3}
La forme quadratique du noyau doit être positive sur les allers-retours (définie positivité)~; pour exclure aussi la manipulation déclenchée par les transactions, la décroissance doit être
convexe et non croissante.
\iqlookfor{Définie positivité~; convexité.}
\end{solution}

\begin{solution}{iq:mx:transient-impact-and-propagator-models:4}
Le premier bloc prend la liquidité déjà présente dans le carnet~; la partie continue s'exécute au rythme de reconstitution du carnet~; le dernier bloc prend une liquidité qui
ne devra jamais être payée par des transactions ultérieures, puisque l'exécution s'achève.
\iqlookfor{Résilience~; effets de bord.}
\end{solution}

\begin{solution}{iq:mx:transient-impact-and-propagator-models:5}
$\tfrac12x^\top\Gamma x$ est en $O(n^2)$~; pour un noyau exponentiel, une somme cumulée le rend en $O(n)$, pour une somme d'exponentielles en $O(nk)$, et pour des noyaux généraux une
convolution par FFT sur une grille régulière est en $O(n\log n)$.
\iqlookfor{Complexité~; astuces récursives ou FFT.}
\end{solution}

\begin{solution}{iq:mx:transient-impact-and-propagator-models:6}
Probablement pas un bogue~: le noyau ajusté n'est pas convexe (ou pas défini positif), et l'optimiseur l'exploite. Il faut corriger le modèle, en ajustant
dans la classe des noyaux convexes ou en projetant l'ajustement, plutôt que l'optimiseur.
\iqlookfor{Manipulation déclenchée par les transactions~; contraintes du modèle.}
\end{solution}
